% !TEX output_directory = ../publica
\documentclass[twoside, 10pt, a4paper]{article}
\usepackage[utf8]{inputenc}
\usepackage[T1]{fontenc}
\usepackage[portuguese]{babel}

% --- FONTE PADRÃO (Estilo Didático Encorpado) ---
\usepackage{helvet} 
\renewcommand{\familydefault}{\sfdefault}

% --- GEOMETRIA E ESPAÇAMENTO ---
\usepackage[top=1.6cm, bottom=1.5cm, left=0.9cm, right=0.9cm, headheight=22pt, headsep=0.4cm]{geometry}
\usepackage{microtype} 
\usepackage{setspace}
\setstretch{1.0}
\usepackage{parskip}
\setlength{\parskip}{2pt}
\setlength{\parindent}{0pt}

\newlength{\espacoPosTitulo}\setlength{\espacoPosTitulo}{0.3cm}
\newlength{\espacoPreQuestao}\setlength{\espacoPreQuestao}{0.1cm}
\newlength{\espacoQuestao}\setlength{\espacoQuestao}{0.2cm}
\newlength{\espacoAlt}\setlength{\espacoAlt}{0.2cm}

\usepackage{enumitem}
\setlist{itemsep=0.4cm, topsep=0.1cm, parsep=0pt, partopsep=0pt}
\newcommand{\larguraTikz}{0.85\linewidth} 

% --- MATEMÁTICA E CORES ---
\usepackage{amsmath, amsfonts, amssymb, nccmath, mathastext}
\usepackage{xcolor}
\definecolor{indigo}{HTML}{4F46E5}
\definecolor{CorTitUm}{HTML}{FF6F61}
\definecolor{CorTitDois}{HTML}{FF6F61}
\definecolor{CorTitTres}{HTML}{658894}
\definecolor{CorLinha}{HTML}{B0B0B0} 
\definecolor{metal}{HTML}{7F8C8D}
\definecolor{vidro}{HTML}{3498DB}
\definecolor{ouro}{HTML}{F1C40F}
\definecolor{concreto}{HTML}{95A5A6}
\definecolor{madeira}{HTML}{D35400}
\definecolor{CinzaEscuro}{HTML}{4A4A4A} 
\definecolor{CinzaMedio}{HTML}{848484} 
\definecolor{Cinzablack}{HTML}{353535} 

% --- MINI-CABEÇALHO E RODAPÉ AUTORAL (TWOSIDE) ---
\usepackage{fancyhdr}
\pagestyle{fancy}
\fancyhf{} 
\renewcommand{\headrulewidth}{0.3pt} 
\renewcommand{\footrulewidth}{0.3pt}
\fancyhead[LE]{\small\sffamily\bfseries\color{gray!75} GEOMETRIA ESPACIAL E VISTAS ORTOGONAIS}
\fancyhead[RO]{\small\sffamily\color{gray!75} \texttt{www.profraphaelpaulino.com.br}}
\fancyfoot[LE]{\small \textbf{\thepage}}
\fancyfoot[RO]{\small \textbf{\thepage}}

% --- CAIXAS DE DESTAQUE ---
\usepackage[most]{tcolorbox}
\newtcolorbox{boxresponda}{colback=white, colframe=gray!45, boxrule=0pt, leftrule=1.7pt, sharp corners, enlarge left by=0.4cm, width=\linewidth-0.4cm, left=8pt, right=0pt, top=2pt, bottom=2pt, fontupper=\small}
\definecolor{CorDicaBorda}{HTML}{17c1be} 
\definecolor{CorDicaFundo}{HTML}{f3fbfb} 
\newtcolorbox{boxdica}{colback=CorDicaFundo, colframe=CorDicaBorda, boxrule=0pt, leftrule=2.5pt, sharp corners, width=\linewidth, left=8pt, right=8pt, top=6pt, bottom=6pt, halign=center, fontupper=\small}

% --- TÍTULOS ---
\usepackage{titlesec}
\titleformat{\section}{\color{black}\large\bfseries}{\thesection}{0em}{\vspace{0.1cm}\titlerule[0.8pt]}
\titlespacing*{\section}{0pt}{10pt}{6pt} 
\titleformat{\subsubsection}[runin]{\color{black}\bfseries}{}{0em}{}
\titlespacing*{\subsubsection}{0pt}{0pt}{0.15cm}

% --- COLUNAS E GRÁFICOS ---
\usepackage{multicol}
\setlength{\columnsep}{1cm}
\setlength{\columnseprule}{0.5pt}
\def\columnseprulecolor{\color{CorLinha}}
\usepackage{adjustbox, tikz, pgfplots, circuitikz}
\usetikzlibrary{shadows, shapes.misc, positioning, calc}
\pgfplotsset{compat=1.18}

\begin{document}
\thispagestyle{empty}
\color{Cinzablack}
\vspace*{-1.8cm}

% --- CABEÇALHO PRINCIPAL AUTORAL (Apenas 1ª página) ---
\noindent
\begin{minipage}{\textwidth}
    \fontfamily{phv}\selectfont
    \noindent
    \begin{minipage}[t]{0.70\linewidth}
        {\Large \bfseries \MakeUppercase{Caderno de Atividades Suplementares}}\\[0.1cm]
        {\large \textmd{\textcolor{CinzaEscuro}{Unidade: GEOMETRIA ESPACIAL E VISTAS ORTOGONAIS}}}
    \end{minipage}%
    \begin{minipage}[t]{0.30\linewidth}
        \raggedleft
        \small \textbf{Prof. Raphael Paulino}\\
        \textsf{\small \textcolor{indigo}{\texttt{profraphaelpaulino.com.br}}}
    \end{minipage}
    
    \vspace{0.15cm}
    \noindent\rule{\linewidth}{1.2pt}
    \vspace{-0.1cm}
    \footnotesize\textsf{\textbf{ÁREA:} MATEMÁTICA / GEOMETRIA \hfill \textbf{NÍVEL:} EF II (9º ANO)}
    \vspace{0.1cm}
    \noindent\rule{\linewidth}{0.4pt}
\end{minipage}

\par\vspace{0.3cm} 
\raggedcolumns
\begin{multicols*}{2}
\vspace{0.1cm}

\subsection*{Capítulo 1: Noções Básicas de Geometria Espacial}
\vspace{-0.2cm}\noindent\rule{\linewidth}{1.5pt} 
\par\vspace{\espacoPosTitulo}
\subsubsection*{1.} Um designer de embalagens está projetando uma caixa especial convexa que possui exatamente \textbf{14} faces, sendo \textbf{8} faces triangulares e \textbf{6} faces octogonais. Para encomendar as conexões metálicas que unirão as arestas nos vértices, ele precisa determinar o número total de vértices dessa embalagem. Modele algebricamente o cálculo do número de arestas e, em seguida, aplique a Relação de Euler para determinar o número de vértices.
\par\vspace{\espacoQuestao}

\noindent\textcolor{CorLinha}{\rule{\linewidth}{0.5pt}} 
\par\vspace{\espacoPreQuestao}
\subsubsection*{2.} Na montagem de um galpão industrial, um engenheiro define o piso nivelado como um plano \textbf{alpha} e um cabo de aço tensionado no teto como uma reta \textbf{r}, estritamente paralela ao plano \textbf{alpha}. Uma máquina, representada por um ponto \textbf{P}, repousa no piso (\textbf{P} pertence a \textbf{alpha}). A partir da máquina \textbf{P}, quantas canaletas retilíneas, totalmente contidas no piso \textbf{alpha}, podem ser instaladas de modo que sejam perfeitamente paralelas ao cabo de aço \textbf{r}? Justifique sua resposta empregando os postulados da geometria de posição.
\par\vspace{\espacoQuestao}

\noindent\textcolor{CorLinha}{\rule{\linewidth}{0.5pt}}
\par\vspace{\espacoPreQuestao}
\subsubsection*{3.} Durante uma aula de laboratório sobre estruturas geométricas, um grupo de alunos discutiu as posições relativas entre as hastes de suporte de uma ponte. O líder do grupo anotou a seguinte conclusão no quadro: "Se pegarmos duas hastes metálicas que não se cruzam (não têm ponto em comum) em nenhum lugar do espaço, podemos garantir matematicamente que essas duas hastes são paralelas."
\begin{boxresponda}
\textbf{Responda:} O raciocínio registrado pelo aluno está correto? Diagnostique a falha conceitual (se houver) explicando qual posição relativa específica no espaço tridimensional ele ignorou, e dê um exemplo no cotidiano que prove o erro de sua afirmação.
\end{boxresponda}
\par\vspace{\espacoQuestao}

\noindent\textcolor{CorLinha}{\rule{\linewidth}{0.5pt}} 
\par\vspace{\espacoPreQuestao}
\subsubsection*{4.} O bloco de vidro a seguir possui suportes metálicos nas suas bordas. Considere as arestas prolongadas como retas e os vértices como pontos no espaço tridimensional.
\begin{center}
% ==========================================
% CONTROLE INDIVIDUAL DESTA IMAGEM:
% 1. CORTAR BORDAS: Mude os zeros do trim (Esquerda, Baixo, Direita, Topo). Ex: trim=1cm 0cm 1cm 0cm
% 2. TAMANHO: Para encolher só esta imagem, troque \larguraTikz por um valor (Ex: 0.5\linewidth)
% ==========================================
\begin{adjustbox}{trim=0cm 0cm 0cm 0cm, clip, max width=\larguraTikz}
\begin{tikzpicture}
% LAYER 1: Sombra
% --- APLICANDO DROP SHADOW ---
\fill[gray!30, drop shadow={opacity=0.3, shadow xshift=3pt, shadow yshift=-3pt}, rounded corners=2pt] (0.5, -0.5) rectangle (3.5, 2.5);
% LAYER 2: Prisma/Bloco de Vidro
\draw[thick, metal, fill=vidro!15, rounded corners=1pt] (0,0) -- (3,0) -- (3,3) -- (0,3) -- cycle;
\draw[thick, metal, fill=vidro!10, rounded corners=1pt] (3,0) -- (4,1) -- (4,4) -- (3,3) -- cycle;
\draw[thick, metal, fill=vidro!20, rounded corners=1pt] (0,3) -- (1,4) -- (4,4) -- (3,3) -- cycle;
% Arestas ocultas
\draw[thick, metal, dashed, rounded corners=1pt] (0,0) -- (1,1) -- (4,1);
\draw[thick, metal, dashed, rounded corners=1pt] (1,1) -- (1,4);
% Vértices e Rótulos com máscara (fill=white)
\node[fill=white, inner sep=1pt, text=black] at (-0.2, -0.2) {\textbf{A}};
\node[fill=white, inner sep=1pt, text=black] at (3.2, -0.2) {\textbf{B}};
\node[fill=white, inner sep=1pt, text=black] at (4.2, 0.8) {\textbf{C}};
\node[fill=white, inner sep=1pt, text=black] at (1.2, 0.8) {\textbf{D}};
\node[fill=white, inner sep=1pt, text=black] at (-0.2, 3.2) {\textbf{E}};
\node[fill=white, inner sep=1pt, text=black] at (3.2, 3.2) {\textbf{F}};
\node[fill=white, inner sep=1pt, text=black] at (4.2, 4.2) {\textbf{G}};
\node[fill=white, inner sep=1pt, text=black] at (1.2, 4.2) {\textbf{H}};
\end{tikzpicture}
\end{adjustbox}
\end{center}
Julgue as afirmativas abaixo acerca das posições relativas presentes na estrutura. Indique (V) para verdadeiro e (F) para falso, justificando rigorosamente os itens falsos.
\begin{enumerate}[label=\alph*), itemsep=\espacoAlt]
\item As retas determinadas pelas arestas \textbf{AB} e \textbf{CG} são ortogonais, logo, elas são concorrentes em um ponto do espaço.
\item A reta que contém a aresta \textbf{EH} e a reta que contém a aresta \textbf{BC} são exemplos clássicos de retas reversas, pois não são paralelas e não possuem ponto de interseção.
\item Os pontos \textbf{A}, \textbf{F}, \textbf{G} e \textbf{D} são estritamente coplanares, ou seja, pertencem a um mesmo plano que "corta" diagonalmente o bloco.
\end{enumerate}
\par\vspace{\espacoQuestao}

\noindent\textcolor{CorLinha}{\rule{\linewidth}{0.5pt}} 
\par\vspace{\espacoPreQuestao}
\subsubsection*{5.} Uma fabricante de bolas de couro estuda um novo design costurado formado exclusivamente por \textbf{x} polígonos hexagonais e \textbf{y} polígonos pentagonais, com todos os vértices perfurando a superfície esférica sem "buracos". Construa a expressão algébrica literal fechada que determina o número total de arestas \textbf{A} dessa estrutura em função das variáveis \textbf{x} e \textbf{y}.
\par\vspace{\espacoQuestao}

\noindent\textcolor{CorLinha}{\rule{\linewidth}{0.5pt}} 
\par\vspace{\espacoPreQuestao}
\subsubsection*{6.} (ENEM - Adaptada) Um lapidador recebeu um cristal bruto e, utilizando cortes precisos a laser, esculpiu a gema no formato de um octaedro regular, apoiado sobre uma base expositora. A estrutura está ancorada por suportes metálicos ao longo de certas arestas.
\begin{center}
% ==========================================
% CONTROLE INDIVIDUAL DESTA IMAGEM:
% 1. CORTAR BORDAS: Mude os zeros do trim (Esquerda, Baixo, Direita, Topo). Ex: trim=1cm 0cm 1cm 0cm
% 2. TAMANHO: Para encolher só esta imagem, troque \larguraTikz por um valor (Ex: 0.5\linewidth)
% ==========================================
\begin{adjustbox}{trim=0cm 0cm 0cm 0cm, clip, max width=\larguraTikz}
\begin{tikzpicture}
% Fundo expositor
% --- APLICANDO DROP SHADOW ---
\fill[gray!10, rounded corners=4pt] (-2.5, -3.5) rectangle (2.5, 3.5);
\fill[gray!80!black, drop shadow={opacity=0.4, shadow xshift=2pt, shadow yshift=-2pt}] (-2, -3) -- (2, -3) -- (2.5, -2.5) -- (-1.5, -2.5) -- cycle;
% Octaedro (Metade Superior)
\draw[thick, metal, fill=vidro!30, opacity=0.9] (0,2.5) -- (-1.5,0) -- (1, -0.5) -- cycle;
\draw[thick, metal, fill=vidro!50, opacity=0.8] (0,2.5) -- (1, -0.5) -- (2, 0.5) -- cycle;
% Octaedro (Metade Inferior)
\draw[thick, metal, fill=vidro!70, opacity=0.9] (-1.5,0) -- (1, -0.5) -- (0, -2) -- cycle;
\draw[thick, metal, fill=vidro!90, opacity=0.8] (1, -0.5) -- (2, 0.5) -- (0, -2) -- cycle;
% Linhas Ocultas
\draw[dashed, metal, thick] (-1.5,0) -- (-0.5, 1) -- (2, 0.5);
\draw[dashed, metal, thick] (0, 2.5) -- (-0.5, 1);
\draw[dashed, metal, thick] (0, -2) -- (-0.5, 1);
% Rótulos
\node[fill=white, inner sep=1pt, text=black] at (0, 2.8) {\textbf{Topo} (\textbf{T})};
\node[fill=white, inner sep=1pt, text=black] at (0, -2.3) {\textbf{Base} (\textbf{B})};
\node[fill=white, inner sep=1pt, text=black] at (-1.8, 0) {\textbf{V}};
\node[fill=white, inner sep=1pt, text=black] at (2.3, 0.5) {\textbf{W}};
\end{tikzpicture}
\end{adjustbox}
\end{center}
Em relação às propriedades espaciais desse cristal e considerando a reta que liga o \textbf{Topo} (\textbf{T}) à \textbf{Base} (\textbf{B}) e a reta que contém a aresta oculta \textbf{VW} (passando por dentro do equador da joia), é correto afirmar que essas retas são:
\begin{enumerate}[label=\alph*), itemsep=\espacoAlt]
\item paralelas e coplanares.
\item concorrentes no centroide do octaedro.
\item reversas, não coplanares e mantêm a mesma distância ao longo do espaço.
\item ortogonais, mas não se interceptam (são reversas).
\item perpendiculares, cruzando-se internamente em um ângulo reto.
\end{enumerate}
\par\vspace{\espacoQuestao}

\noindent\textcolor{CorLinha}{\rule{\linewidth}{0.5pt}} 
\par\vspace{\espacoPreQuestao}
\subsubsection*{7.} (Estilo VUNESP) Em um projeto arquitetônico, um museu foi desenhado como um grande poliedro convexo maciço. Sabe-se que esse poliedro possui \textbf{20} arestas e apenas dois tipos de faces: \textbf{f} faces quadrangulares e \textbf{t} faces triangulares. Além disso, a contagem de vértices revelou que a estrutura tem exatamente \textbf{9} vértices. O engenheiro estagiário tentou calcular quantas faces triangulares o projeto possui. Prove matematicamente qual é o valor correto de \textbf{t} e apresente a justificativa algébrica completa que refuta a ideia de que esse poliedro poderia ter \textbf{10} faces quadrangulares.
\par\vspace{\espacoQuestao}

\noindent\textcolor{CorLinha}{\rule{\linewidth}{0.5pt}} 
\par\vspace{\espacoPreQuestao}
\subsubsection*{8.} A tesoura de telhado é uma estrutura treliçada de madeira essencial para sustentar a cobertura de residências. O esquema estrutural abaixo detalha a disposição das vigas mestras.
\begin{center}
% ==========================================
% CONTROLE INDIVIDUAL DESTA IMAGEM:
% 1. CORTAR BORDAS: Mude os zeros do trim (Esquerda, Baixo, Direita, Topo). Ex: trim=1cm 0cm 1cm 0cm
% 2. TAMANHO: Para encolher só esta imagem, troque \larguraTikz por um valor (Ex: 0.5\linewidth)
% ==========================================
\begin{adjustbox}{trim=0cm 0cm 0cm 0cm, clip, max width=\larguraTikz}
\begin{tikzpicture}
\fill[gray!5] (-4,-1) rectangle (4, 3.5);
% Vigas da Tesoura (Traços Grossos com Sombra)
% --- APLICANDO DROP SHADOW ---
\draw[line width=4pt, madeira!80!black, drop shadow={opacity=0.3, shadow xshift=2pt, shadow yshift=-2pt}] (-3,0) -- (3,0); 
\draw[line width=4pt, madeira!80!black] (-3,0) -- (0,2.5) -- (3,0); 
\draw[line width=3pt, madeira!60!black] (0,0) -- (0,2.5); 
\draw[line width=3pt, madeira!60!black] (-1.5, 0) -- (-1.5, 1.25); 
\draw[line width=3pt, madeira!60!black] (1.5, 0) -- (1.5, 1.25); 
\draw[line width=3pt, madeira!60!black] (-3,0) -- (-1.5, 1.25); 
\draw[line width=3pt, madeira!60!black] (0,0) -- (-1.5, 1.25);
\draw[line width=3pt, madeira!60!black] (0,0) -- (1.5, 1.25);
\draw[line width=3pt, madeira!60!black] (3,0) -- (1.5, 1.25);
% Nós / Rótulos
\fill[black] (-3,0) circle (2pt) node[below, fill=white, inner sep=1pt] {\textbf{A}};
\fill[black] (3,0) circle (2pt) node[below, fill=white, inner sep=1pt] {\textbf{B}};
\fill[black] (0,2.5) circle (2pt) node[above, fill=white, inner sep=1pt] {\textbf{C}};
\fill[black] (0,0) circle (2pt) node[below, fill=white, inner sep=1pt] {\textbf{D}};
\fill[black] (-1.5,1.25) circle (2pt) node[above left, fill=white, inner sep=1pt] {\textbf{E}};
\end{tikzpicture}
\end{adjustbox}
\end{center}
Supondo que a reta de suporte da viga horizontal inferior seja perpendicular ao plano das paredes (não desenhadas), defina a posição geométrica relativa (concorrentes, paralelas ou reversas) entre a reta suporte da viga \textbf{CD} (pendural central) e a reta que contém a viga \textbf{AE}, baseando-se unicamente nas intersecções visíveis neste plano bidimensional de corte da estrutura.
\par\vspace{\espacoQuestao}

\subsection*{Capítulo 2: Vistas Ortogonais}
\vspace{-0.2cm}\noindent\rule{\linewidth}{1.5pt} 
\par\vspace{\espacoPosTitulo}
\subsubsection*{9.} Um aluno de marcenaria montou a escultura isométrica abaixo utilizando pequenos cubos de madeira maciça idênticos, colando-os face a face.
\begin{center}
% ==========================================
% CONTROLE INDIVIDUAL DESTA IMAGEM:
% 1. CORTAR BORDAS: Mude os zeros do trim (Esquerda, Baixo, Direita, Topo). Ex: trim=1cm 0cm 1cm 0cm
% 2. TAMANHO: Para encolher só esta imagem, troque \larguraTikz por um valor (Ex: 0.5\linewidth)
% ==========================================
\begin{adjustbox}{trim=0cm 0cm 0cm 0cm, clip, max width=\larguraTikz}
\begin{tikzpicture}[x=(210:0.5cm), y=(-30:0.5cm), z=(90:0.5cm)]
\definecolor{cubo}{HTML}{D4AC0D}
% --- APLICANDO ARREDONDAMENTO DE CANTOS E PREENCHIMENTO ---
\tikzstyle{face_top} = [fill=cubo!40, draw=black, thick, rounded corners=0.5pt]
\tikzstyle{face_left} = [fill=cubo!80, draw=black, thick, rounded corners=0.5pt]
\tikzstyle{face_right} = [fill=cubo!100, draw=black, thick, rounded corners=0.5pt]

\def\cubo#1#2#3{
  \draw[face_top] (#1,#2,#3+1) -- (#1+1,#2,#3+1) -- (#1+1,#2+1,#3+1) -- (#1,#2+1,#3+1) -- cycle;
  \draw[face_left] (#1+1,#2,#3) -- (#1+1,#2+1,#3) -- (#1+1,#2+1,#3+1) -- (#1+1,#2,#3+1) -- cycle;
  \draw[face_right] (#1,#2+1,#3) -- (#1+1,#2+1,#3) -- (#1+1,#2+1,#3+1) -- (#1,#2+1,#3+1) -- cycle;
}

% Base (Z=0)
\cubo{0}{0}{0}
\cubo{0}{1}{0}
\cubo{0}{2}{0}
\cubo{1}{0}{0}
% Nivel 1 (Z=1)
\cubo{0}{1}{1}
\cubo{0}{2}{1}
% Nivel 2 (Z=2)
\cubo{0}{2}{2}

\draw[->, ultra thick, red!80!black, drop shadow={opacity=0.3}] (2, -1, 0) -- (1, 0, 0) node[midway, below right, fill=white, inner sep=1pt] {\textbf{Frontal}};

\end{tikzpicture}
\end{adjustbox}
\end{center}
Desenhe o esboço ou decreva detalhadamente (indicando colunas e alturas) como ficarão a Vista Superior (olhando de cima) e a Vista Lateral Esquerda desta montagem tridimensional.
\par\vspace{\espacoQuestao}

\noindent\textcolor{CorLinha}{\rule{\linewidth}{0.5pt}}
\par\vspace{\espacoPreQuestao}
\subsubsection*{10.} Ao receber a mesma peça isométrica da questão anterior para análise, um torneiro mecânico desenhou a Vista Superior formando apenas um retângulo contínuo de proporção \textbf{3} por \textbf{1}, sem nenhuma linha interna dividindo os quadrados.
\begin{boxresponda}
\textbf{Responda:} Qual erro técnico de representação o torneiro cometeu ao ignorar as linhas internas da Vista Superior? O que as linhas contínuas internas em uma vista ortogonal revelam sobre a topografia (alturas/profundidades) do objeto tridimensional que ele falhou em representar?
\end{boxresponda}
\par\vspace{\espacoQuestao}

\noindent\textcolor{CorLinha}{\rule{\linewidth}{0.5pt}} 
\par\vspace{\espacoPreQuestao}
\subsubsection*{11.} Considere um cilindro maciço de aço posicionado em pé sobre uma bancada plana. Julgue as afirmativas abaixo, relacionadas à representação técnica desse cilindro, marcando V ou F e justificando os possíveis erros das afirmativas.
\begin{enumerate}[label=\Roman*., itemsep=\espacoAlt]
\item A Vista Superior de um cilindro em pé é perfeitamente representada por uma circunferência geométrica.
\item A Vista Frontal desse cilindro apresentará linhas curvas nas laterais, acompanhando o raio do objeto.
\item A Vista Lateral Direita e a Vista Frontal deste cilindro em pé gerarão projeções retangulares absolutamente idênticas na folha de projeto.
\end{enumerate}
\par\vspace{\espacoQuestao}

\noindent\textcolor{CorLinha}{\rule{\linewidth}{0.5pt}} 
\par\vspace{\espacoPreQuestao}
\subsubsection*{12.} (ENEM - Adaptada) O sol a pino ao meio-dia gera sombras que, geometricamente, se comportam como Vistas Superiores exatas dos objetos iluminados em relação ao chão plano. Uma prefeitura instalou no parque uma escultura abstrata metálica com o formato de um tronco de cone circular reto vazado (como um copo de ponta-cabeça sem o fundo superior e sem a base). Desconsiderando a espessura da chapa de metal, a sombra projetada por essa estrutura no piso horizontal perfeitamente plano ao meio-dia solar será identificada visualmente como:
\begin{enumerate}[label=\alph*), itemsep=\espacoAlt]
\item um único círculo maciço e escurecido.
\item dois círculos concêntricos separados por uma região preenchida (coroa circular).
\item um trapézio isósceles com bases retas e laterais inclinadas.
\item uma elipse alongada no sentido da iluminação.
\item apenas a borda de uma circunferência, sem espessura, já que a peça é vazada nos dois extremos.
\end{enumerate}
\par\vspace{\espacoQuestao}

\noindent\textcolor{CorLinha}{\rule{\linewidth}{0.5pt}} 
\par\vspace{\espacoPreQuestao}
\subsubsection*{13.} (ENEM - Adaptada) Um estúdio de design projetou um pódio maciço de premiação esculpido em um único bloco geométrico. O projeto isométrico tridimensional da estrutura está detalhado abaixo.
\begin{center}
% ==========================================
% CONTROLE INDIVIDUAL DESTA IMAGEM:
% 1. CORTAR BORDAS: Mude os zeros do trim (Esquerda, Baixo, Direita, Topo). Ex: trim=1cm 0cm 1cm 0cm
% 2. TAMANHO: Para encolher só esta imagem, troque \larguraTikz por um valor (Ex: 0.5\linewidth)
% ==========================================
\begin{adjustbox}{trim=0cm 0cm 0cm 0cm, clip, max width=\larguraTikz}
\begin{tikzpicture}[x=(210:0.5cm), y=(-30:0.5cm), z=(90:0.5cm)]
% --- APLICANDO DETALHES DE CONCRETO E PREENCHIMENTO ---
\tikzstyle{face_top} = [fill=concreto!40, draw=CinzaEscuro, thick, rounded corners=0.5pt]
\tikzstyle{face_left} = [fill=concreto!70, draw=CinzaEscuro, thick, rounded corners=0.5pt]
\tikzstyle{face_right} = [fill=concreto!100, draw=CinzaEscuro, thick, rounded corners=0.5pt]

\def\cubo#1#2#3{
  \draw[face_top] (#1,#2,#3+1) -- (#1+1,#2,#3+1) -- (#1+1,#2+1,#3+1) -- (#1,#2+1,#3+1) -- cycle;
  \draw[face_left] (#1+1,#2,#3) -- (#1+1,#2+1,#3) -- (#1+1,#2+1,#3+1) -- (#1+1,#2,#3+1) -- cycle;
  \draw[face_right] (#1,#2+1,#3) -- (#1+1,#2+1,#3) -- (#1+1,#2+1,#3+1) -- (#1,#2+1,#3+1) -- cycle;
}

% Base nivel 0
\cubo{0}{0}{0} \cubo{0}{1}{0} \cubo{0}{2}{0}
\cubo{1}{0}{0} \cubo{1}{1}{0} \cubo{1}{2}{0}
% Nivel 1 (degrau do meio e alto)
\cubo{0}{1}{1} \cubo{0}{2}{1}
\cubo{1}{1}{1} \cubo{1}{2}{1}
% Nivel 2 (degrau mais alto)
\cubo{0}{2}{2}
\cubo{1}{2}{2}

% Setas de observação
% --- APLICANDO SOMBRAS NAS SETAS ---
\draw[->, ultra thick, ouro, drop shadow={opacity=0.3}] (2.5, -1, 0) -- (1.5, 0, 0) node[midway, below right, fill=white, inner sep=2pt] {\textbf{Vista Frontal}};
\draw[->, ultra thick, vidro, drop shadow={opacity=0.3}] (3, 3, 0) -- (1.5, 2, 0) node[midway, above right, fill=white, inner sep=2pt] {\textbf{Vista Lateral}};
\end{tikzpicture}
\end{adjustbox}
\end{center}
Um marceneiro precisa fabricar o molde bidimensional exato da Vista Frontal dessa peça. Considerando a seta indicativa, assinale a alternativa que descreve o perfil correto dessa vista e justifique matematicamente o erro das alternativas B e C, explicando qual falha de leitura espacial elas representam.
\begin{enumerate}[label=\alph*), itemsep=\espacoAlt]
\item Um retângulo de proporção \textbf{3} por \textbf{1} com dois quadrados empilhados na extremidade direita.
\item Uma escada decrescente da esquerda para a direita, contendo linhas tracejadas.
\item Um quadrado perfeito \textbf{3} por \textbf{3} subdividido em \textbf{9} quadrados menores, todos visíveis.
\item Um contorno em formato de "L" invertido, com \textbf{2} colunas de largura e \textbf{3} de altura.
\item Três colunas retangulares lado a lado, com alturas proporcionais a \textbf{1}, \textbf{2} e \textbf{3} unidades, crescendo da esquerda para a direita.
\end{enumerate}
\par\vspace{\espacoQuestao}

\noindent\textcolor{CorLinha}{\rule{\linewidth}{0.5pt}} 
\par\vspace{\espacoPreQuestao}
\subsubsection*{14.} Durante o desenvolvimento de um jogo 3D, um programador desenhou a Vista Superior de um castelo cuja torre principal é um cilindro perfeito finalizado por um telhado em formato de cone circular reto (o cone está centralizado e tem o mesmo raio da base do cilindro). Ao renderizar a planta baixa, ele desenhou apenas um círculo contínuo com um ponto escuro no centro.
\begin{boxresponda}
\textbf{Responda:} A representação técnica bidimensional (Vista Superior) feita pelo programador está completa para distinguir o cone do cilindro? Diagnostique qual informação geométrica de profundidade ou superfície o "ponto central" tenta traduzir e por que essa vista seria idêntica se a torre fosse apenas uma pirâmide de base circular (cone) diretamente no chão.
\end{boxresponda}
\par\vspace{\espacoQuestao}

\noindent\textcolor{CorLinha}{\rule{\linewidth}{0.5pt}} 
\par\vspace{\espacoPreQuestao}
\subsubsection*{15.} Um bloco maciço foi usinado de modo que sua Vista Frontal é um triângulo isósceles e sua Vista Superior é um retângulo contínuo, com uma linha reta dividindo o retângulo exatamente ao meio no sentido do comprimento. A partir da engenharia reversa dessas vistas bidimensionais, deduza e escreva o nome geométrico tridimensional (sólido) mais provável para essa peça.
\par\vspace{\espacoQuestao}
\noindent\textcolor{CorLinha}{\rule{\linewidth}{0.5pt}} 
\par\vspace{\espacoPreQuestao}
\subsubsection*{16.} Analise as afirmativas a seguir sobre a rotação de objetos no espaço e a alteração de suas vistas ortogonais. Marque (V) para verdadeiro e (F) para falso. Apresente uma justificativa analítica provando o erro das afirmativas falsas.
\begin{enumerate}[label=\Roman*., itemsep=\espacoAlt]
\item Se rotacionarmos um cubo perfeito em \textbf{45 graus} em torno do seu eixo vertical (\textbf{z}), sua Vista Superior deixará de ser um quadrado e passará a ser representada por um losango ou dois retângulos adjacentes.
\item A Vista Frontal de uma esfera maciça é a única representação ortogonal que não sofre nenhuma alteração geométrica, independentemente do eixo ou do ângulo em que a esfera seja rotacionada.
\item Quando uma pirâmide regular de base quadrada é tombada de lado (apoiando uma de suas faces triangulares no chão plano), sua Vista Superior continua apresentando um contorno perfeitamente quadrado.
\end{enumerate}
\par\vspace{\espacoQuestao}

\subsection*{Capítulo 3: Representações Técnicas}
\vspace{-0.2cm}\noindent\rule{\linewidth}{1.5pt} 
\par\vspace{\espacoPosTitulo}
\subsubsection*{17.} Na leitura de desenho técnico mecânico, a relação entre as vistas dita a topografia da peça. Considere um prisma metálico cujo esboço da Vista Superior mostra dois círculos concêntricos (um dentro do outro, com linhas contínuas). A Vista Frontal desse mesmo prisma é um retângulo contendo duas linhas internas tracejadas verticais, que cruzam a peça de cima a baixo. Descreva fisicamente o que essas linhas tracejadas representam no objeto tridimensional real.
\begin{boxdica}
Lembre-se: em desenho técnico, linhas visíveis e contínuas indicam arestas expostas ao observador. Linhas tracejadas indicam contornos, furos ou rebaixos que existem na peça, mas estão ocultos pela massa do material daquela perspectiva.
\end{boxdica}
\par\vspace{\espacoQuestao}

\noindent\textcolor{CorLinha}{\rule{\linewidth}{0.5pt}} 
\par\vspace{\espacoPreQuestao}
\subsubsection*{18.} O professor de robótica distribuiu a planta de uma haste estabilizadora. O aluno Pedro analisou a Vista Lateral Direita, que consistia em um retângulo atravessado horizontalmente por uma linha contínua grossa no terço superior. Pedro imediatamente concluiu: "A peça tem um rasgo/furo invisível que a atravessa nesse ponto".
\begin{boxresponda}
\textbf{Responda:} Qual é o erro técnico de leitura cometido por Pedro ao interpretar a linha bidimensional? Qual é o tipo exato de linha que deveria aparecer no desenho para validar a conclusão dele sobre um "furo invisível"?
\end{boxresponda}
\par\vspace{\espacoQuestao}

\noindent\textcolor{CorLinha}{\rule{\linewidth}{0.5pt}} 
\par\vspace{\espacoPreQuestao}
\subsubsection*{19.} Um suporte de prateleira em formato de "U" invertido é fabricado dobrando-se uma chapa de aço maciça. Se olharmos essa peça diretamente por cima (Vista Superior), veremos apenas um retângulo longo e contínuo. Sabendo que o "U" é oco por baixo, explique por que o desenhista é obrigado a inserir linhas tracejadas nessa Vista Superior e determine exatamente quantas linhas tracejadas longitudinais devem ser traçadas.
\par\vspace{\espacoQuestao}

\noindent\textcolor{CorLinha}{\rule{\linewidth}{0.5pt}} 
\par\vspace{\espacoPreQuestao}
\subsubsection*{20.} O componente mecânico de aço ilustrado abaixo possui um furo cilíndrico passante (que atravessa toda a aba inferior).
\begin{center}
% ==========================================
% CONTROLE INDIVIDUAL DESTA IMAGEM:
% 1. CORTAR BORDAS: Mude os zeros do trim (Esquerda, Baixo, Direita, Topo). Ex: trim=1cm 0cm 1cm 0cm
% 2. TAMANHO: Para encolher só esta imagem, troque \larguraTikz por um valor (Ex: 0.5\linewidth)
% ==========================================
\begin{adjustbox}{trim=0cm 0cm 0cm 0cm, clip, max width=\larguraTikz}
\begin{tikzpicture}[x=(210:0.5cm), y=(-30:0.5cm), z=(90:0.5cm)]
% Layer Sombra
% --- APLICANDO DROP SHADOW E CANTOS ARREDONDADOS ---
\fill[gray!20, drop shadow={opacity=0.4, shadow xshift=4pt, shadow yshift=-4pt}, rounded corners=2pt] (1.5,0.5,-0.5) -- (5,0,-0.5) -- (5,2,-0.5) -- (1,2.5,-0.5) -- cycle;

% Aba inferior (fundo)
\draw[fill=metal!60, draw=CinzaEscuro, thick, rounded corners=1pt] (0,0,0) -- (4,0,0) -- (4,2,0) -- (0,2,0) -- cycle;
% Aba superior (vertical)
\draw[fill=metal!40, draw=CinzaEscuro, thick, rounded corners=1pt] (0,0,0) -- (1,0,0) -- (1,0,3) -- (0,0,3) -- cycle;
\draw[fill=metal!90, draw=CinzaEscuro, thick, rounded corners=1pt] (1,0,0) -- (1,2,0) -- (1,2,3) -- (1,0,3) -- cycle;
\draw[fill=metal!30, draw=CinzaEscuro, thick, rounded corners=1pt] (0,0,3) -- (1,0,3) -- (1,2,3) -- (0,2,3) -- cycle;
\draw[fill=metal!80, draw=CinzaEscuro, thick, rounded corners=1pt] (0,2,0) -- (1,2,0) -- (1,2,3) -- (0,2,3) -- cycle;

% Furo na aba inferior
\draw[fill=white, draw=Cinzablack, thick] (2.5,1,0) circle (0.4cm);
\draw[fill=metal!20, draw=none] (2.4, 0.9, -0.1) circle (0.35cm); % Profundidade do furo
\end{tikzpicture}
\end{adjustbox}
\end{center}
Se um engenheiro precisar projetar rigorosamente a Vista Superior desse componente, qual deverá ser a representação geométrica correta das arestas visíveis? Refute as opções incorretas apontando a omissão geométrica de cada uma.
\begin{enumerate}[label=\alph*), itemsep=\espacoAlt]
\item Um retângulo simples com um círculo no centro, sem nenhuma linha divisória interna.
\item Dois retângulos adjacentes (um representando a espessura da aba vertical e outro a base) e um círculo contido no retângulo maior.
\item Três retângulos paralelos e um furo tracejado no retângulo central.
\item Apenas um contorno em "L" contendo um círculo em uma de suas pernas.
\item Um quadrado perfeito cortado por uma reta diagonal, com um círculo posicionado em um dos vértices.
\end{enumerate}
\par\vspace{\espacoQuestao}

\noindent\textcolor{CorLinha}{\rule{\linewidth}{0.5pt}} 
\par\vspace{\espacoPreQuestao}
\subsubsection*{21.} (FUVEST - Adaptada) O diedro de projeção é o sistema normatizado de organização das vistas em um papel. No Brasil, utiliza-se como padrão o sistema de projeção no 1º Diedro. Sobre as regras rígidas de posicionamento de vistas técnicas neste sistema, julgue as afirmativas apresentando a justificativa estrutural.
\begin{enumerate}[label=\Roman*., itemsep=\espacoAlt]
\item A Vista Superior é obrigatoriamente desenhada abaixo da Vista Frontal, pois a projeção reflete a imagem no plano que está posicionado atrás/abaixo da peça.
\item A Vista Lateral Direita deve ser posicionada à esquerda da Vista Frontal.
\item A cotação (medidas) da peça só pode ser inserida na Vista Frontal, sendo proibido indicar comprimentos nas vistas laterais para evitar redundância.
\end{enumerate}
\par\vspace{\espacoQuestao}

\noindent\textcolor{CorLinha}{\rule{\linewidth}{0.5pt}} 
\par\vspace{\espacoPreQuestao}
\subsubsection*{22.} Em uma planta baixa arquitetônica (que é uma aplicação direta da Vista Superior com um corte horizontal), as paredes são desenhadas com duas linhas paralelas. Uma janela em uma parede é representada apagando-se o preenchimento dessas linhas e desenhando-se traços mais finos no meio. Se um vão de porta não possui porta instalada, apenas a abertura livre, como a continuidade das linhas do piso e da parede deve ser matematicamente representada para garantir a fluidez do trânsito no desenho bidimensional?
\par\vspace{\espacoQuestao}

\subsection*{Capítulo 4: Sistema de Medidas}
\vspace{-0.2cm}\noindent\rule{\linewidth}{1.5pt} 
\par\vspace{\espacoPosTitulo}
\subsubsection*{23.} A astronomia utiliza o Sistema Internacional para mensurar escalas que fogem à percepção humana. A distância média entre a Terra e Marte em seu ponto de maior aproximação é de cerca de \textbf{54,6} milhões de quilômetros. Uma sonda espacial viaja a uma velocidade constante de $ \mfrac{3}{10} $ da velocidade da luz (onde a velocidade da luz $c = 3 \cdot 10^8$ m/s). Converta a distância até Marte estritamente para metros em notação científica e modele a equação algébrica que determina o tempo de viagem dessa sonda em segundos.
\par\vspace{\espacoQuestao}

\noindent\textcolor{CorLinha}{\rule{\linewidth}{0.5pt}} 
\par\vspace{\espacoPreQuestao}
\subsubsection*{24.} Ao conectar um pen drive novo de "32 GB" no computador, um aluno percebe que o sistema operacional exibe a capacidade máxima como sendo apenas aproximadamente \textbf{29,8 GB}. Ele reclama na loja alegando que faltam mais de 2 Gigabytes no produto. O vendedor de TI tenta explicar que a culpa não é da fabricante, mas da matemática base computacional versus a matemática comercial.
\begin{boxresponda}
\textbf{Responda:} Qual é a raiz matemática (base numérica) exata que gera essa divergência estrutural de medidas informáticas? Diagnostique o erro de conversão multiplicativa que o cérebro do aluno faz (base 10) versus o que o processador do computador calcula rigorosamente (base 2) na transição de Bytes para Kilobytes e Gigabytes.
\end{boxresponda}
\begin{boxdica}
Atenção à nomenclatura estrutural do SI: o prefixo "Kilo" no comércio significa $10^3$ (\textbf{1000}). No entanto, o sistema binário da máquina agrupa os dados em potências de \textbf{2}. Qual é a potência de base \textbf{2} que mais se aproxima de \textbf{1000}?
\end{boxdica}
\par\vspace{\espacoQuestao}

\noindent\textcolor{CorLinha}{\rule{\linewidth}{0.5pt}} 
\par\vspace{\espacoPreQuestao}
\subsubsection*{25.} A nanotecnologia opera em escalas imperceptíveis ao olho humano. Enquanto o diâmetro médio de um fio de cabelo humano mede cerca de \textbf{80} micrometros (\textbf{80} $\mu$m), um nanotubo de carbono de parede única possui um diâmetro de aproximadamente \textbf{2} nanômetros (\textbf{2} nm). Sabendo que as grandezas estão no Sistema Internacional, modele a equação utilizando notação científica e determine quantos nanotubos de carbono caberiam, perfeitamente alinhados lado a lado, na espessura de um único fio de cabelo. Utilize o formato de fração $ \mfrac{a}{b} $ na armação do seu cálculo analítico.
\par\vspace{\espacoQuestao}

\noindent\textcolor{CorLinha}{\rule{\linewidth}{0.5pt}} 
\par\vspace{\espacoPreQuestao}
\subsubsection*{26.} Um data center de armazenamento em nuvem organiza seus dados em racks físicos modulares, cuja estrutura de uma gaveta (blade) frontal está ilustrada abaixo.
\begin{center}
% ==========================================
% CONTROLE INDIVIDUAL DESTA IMAGEM:
% 1. CORTAR BORDAS: Mude os zeros do trim (Esquerda, Baixo, Direita, Topo). Ex: trim=1cm 0cm 1cm 0cm
% 2. TAMANHO: Para encolher só esta imagem, troque \larguraTikz por um valor (Ex: 0.5\linewidth)
% ==========================================
\begin{adjustbox}{trim=0cm 0cm 0cm 0cm, clip, max width=\larguraTikz}
\begin{tikzpicture}
% Sombra do Rack
% --- APLICANDO DROP SHADOW ---
\fill[gray!20, drop shadow={opacity=0.5, shadow xshift=4pt, shadow yshift=-4pt}, rounded corners=2pt] (0,0) rectangle (6,4);
% Chassi do Rack
\draw[thick, Cinzablack, fill=metal!40, rounded corners=2pt] (0,0) rectangle (6,4);
\draw[thick, CinzaEscuro, fill=Cinzablack, rounded corners=1pt] (0.2,0.2) rectangle (5.8,3.8);

% Discos Rígidos (Leds e Baias)
\foreach \y in {0.4, 1.3, 2.2, 3.1} {
    \draw[thick, CinzaMedio, fill=metal!80, rounded corners=1pt] (0.4, \y) rectangle (5.6, \y+0.7);
    % Slots de HD
    \foreach \x in {0.6, 1.4, 2.2, 3.0} {
        \draw[thick, Cinzablack, fill=CinzaEscuro] (\x, \y+0.1) rectangle (\x+0.6, \y+0.6);
        % Led de atividade
        \fill[CorDicaBorda] (\x+0.1, \y+0.5) circle (1.5pt);
    }
    % Etiqueta de capacidade (Máscara branca para leitura)
    \node[fill=white, inner sep=1.5pt, rounded corners=1pt] at (4.7, \y+0.35) {\scriptsize \textbf{16 TB}};
}
\end{tikzpicture}
\end{adjustbox}
\end{center}
Cada linha da gaveta acima representa um conjunto de processamento. Sabe-se que cada conjunto suporta exatamente \textbf{4} discos rígidos de \textbf{16} Terabytes (\textbf{TB}) cada. Calcule a capacidade máxima de armazenamento bruto dessa gaveta estrutural em \textbf{TB}. Em seguida, prove matematicamente por que a conversão desse volume total para Gigabytes (\textbf{GB}) no rigor do sistema binário da máquina (onde \textbf{1 TB} equivale a \textbf{1024 GB}) resulta em um número sensivelmente maior do que a conversão comercial praticada pelas lojas na base 10.
\par\vspace{\espacoQuestao}

\noindent\textcolor{CorLinha}{\rule{\linewidth}{0.5pt}} 
\par\vspace{\espacoPreQuestao}
\subsubsection*{27.} (ENEM - Adaptada) O fluxo de transferência de dados em redes de computadores é medido em bits por segundo (bps), enquanto o armazenamento de arquivos é medido em Bytes (B). Sabe-se na arquitetura de computadores que \textbf{1} Byte é composto exatamente por um pacote de \textbf{8} bits. Um estudante contratou um plano de internet de fibra óptica com velocidade real e constante de download de \textbf{400} Megabits por segundo (Mbps). Ele iniciou o download de um jogo em rede que ocupa um espaço de armazenamento maciço de \textbf{90} Gigabytes (GB). Desconsiderando instabilidades na rede e considerando os prefixos do Sistema Internacional na base 10 para facilitar o cálculo do tempo estimado (onde Mega = $10^6$ e Giga = $10^9$), o tempo teórico mínimo, em minutos, para a conclusão desse download será de:
\begin{enumerate}[label=\alph*), itemsep=\espacoAlt]
\item \textbf{30} minutos.
\item \textbf{3} minutos.
\item \textbf{240} minutos.
\item \textbf{24} minutos.
\item \textbf{45} minutos.
\end{enumerate}
Apresente os cálculos de conversão e justifique matematicamente o erro da alternativa C, demonstrando qual conversão estrutural de unidades o candidato esqueceu de aplicar no cruzamento entre fluxo e armazenamento.
\par\vspace{\espacoQuestao}

\noindent\textcolor{CorLinha}{\rule{\linewidth}{0.5pt}} 
\par\vspace{\espacoPreQuestao}
\subsubsection*{28.} Em um projeto escolar de engenharia para criar o contrapiso de uma maquete, um grupo precisa encher um molde tridimensional de concreto. O volume interno calculado foi de \textbf{0,5} m³. O aluno responsável pela dosagem do material fez a seguinte anotação no relatório do grupo: "Como sabemos perfeitamente que a conversão direta de metros para centímetros é multiplicar por \textbf{100}, então o volume da nossa maquete será de $0,5 \times 100 = 50 \text{ cm}^3$. Logo, apenas um copo americano de cimento será suficiente para preencher o piso."
\begin{boxresponda}
\textbf{Responda:} Qual erro dimensional o estudante cometeu ao aplicar a conversão linear de base decimal direta em uma grandeza espacial de volume? Demonstre passo a passo a dedução geométrica das arestas de um cubo de \textbf{1} m para provar qual é a taxa de conversão correta de metros cúbicos para centímetros cúbicos e apresente o volume verdadeiro da maquete.
\end{boxresponda}
\par\vspace{\espacoQuestao}

\noindent\textcolor{CorLinha}{\rule{\linewidth}{0.5pt}} 
\par\vspace{\espacoPreQuestao}
\subsubsection*{29.} A astronomia contemporânea utiliza as propriedades de potenciação do Sistema Internacional para mensurar vazios que fogem à percepção humana. A distância média da Terra à estrela Alpha Centauri é de aproximadamente \textbf{4,3} anos-luz. Sabendo que a luz viaja no vácuo a uma velocidade extrema de $3 \cdot 10^5 \text{ km/s}$ e assumindo, para fins de aproximação estrutural, que um ano terrestre possui cerca de $3 \cdot 10^7$ segundos, determine o valor algébrico exato dessa distância em quilômetros. O resultado final deve ser escrito obrigatoriamente utilizando as regras de notação científica padronizada.
\par\vspace{\espacoQuestao}

\noindent\textcolor{CorLinha}{\rule{\linewidth}{0.5pt}} 
\par\vspace{\espacoPreQuestao}
\subsubsection*{30.} O avanço do poder de processamento global depende puramente da miniaturização geométrica. O diagrama a seguir representa um recorte da vista ortogonal de um "die" (núcleo de silício) de um processador.
\begin{center}
% ==========================================
% CONTROLE INDIVIDUAL DESTA IMAGEM:
% 1. CORTAR BORDAS: Mude os zeros do trim (Esquerda, Baixo, Direita, Topo). Ex: trim=1cm 0cm 1cm 0cm
% 2. TAMANHO: Para encolher só esta imagem, troque \larguraTikz por um valor (Ex: 0.5\linewidth)
% ==========================================
\begin{adjustbox}{trim=0cm 0cm 0cm 0cm, clip, max width=\larguraTikz}
\begin{tikzpicture}
% Fundo do Chip de Silício
% --- APLICANDO DROP SHADOW ---
\fill[Cinzablack, drop shadow={opacity=0.4, shadow xshift=3pt, shadow yshift=-3pt}, rounded corners=2pt] (0,0) rectangle (4,4);
\draw[ouro, line width=1.5pt] (0.2,0.2) rectangle (3.8,3.8);

% Pinos e trilhas externas
\foreach \x in {0.6, 1.3, 2.0, 2.7, 3.4} {
    \draw[ouro, thick] (\x, 0) -- (\x, -0.3);
    \draw[ouro, thick] (\x, 4) -- (\x, 4.3);
    \draw[ouro, thick] (0, \x) -- (-0.3, \x);
    \draw[ouro, thick] (4, \x) -- (4.3, \x);
}

% Destaque microscópico (Callout de Lente)
\draw[thick, vidro, fill=CinzaMedio!40, rounded corners=1pt] (1.2, 1.2) rectangle (2.8, 2.8);
\draw[thick, vidro, dashed] (2.8, 2.8) -- (4.5, 4.5);
\draw[thick, vidro, dashed] (2.8, 1.2) -- (4.5, -0.5);

% Janela da Lente (Transistor)
\fill[metal!20, draw=vidro, thick, drop shadow={opacity=0.3}] (4.5, -0.5) rectangle (7.5, 4.5);
\fill[ouro!80!black] (5.5, 1) rectangle (6.5, 3.5); % Gate do Transistor
\draw[<->, thick, Cinzablack] (5.5, 0.5) -- (6.5, 0.5);
% Máscara branca no nó para proteção das retas de fundo
\node[fill=white, inner sep=2pt] at (6, 0.5) {\textbf{14 nm}};
\end{tikzpicture}
\end{adjustbox}
\end{center}
Nesta arquitetura clássica, a largura de controle do fluxo de elétrons indicada pela cota no destaque determina a escala nominal da tecnologia. A fabricante planeja migrar dessa arquitetura atual de \textbf{14} nm para a novíssima litografia de \textbf{2} nm. Considerando que a área total da pastilha do processador permanecerá estritamente a mesma, explique algebricamente por que a redução unidimensional (comprimento) do transistor não gera um ganho linear de espaço. 
\begin{boxdica}
Modelagem mental: Para comparar as densidades, pense no espaço do transistor como sendo um quadrado simples de lado \textbf{L}. A área microscópica ocupada por ele na placa será $L^2$. A proporção bruta de ganho de espaço (quantas peças a mais cabem) é obtida dividindo a área espacial da geração antiga pela nova.
\end{boxdica}
Calcule geometricamente quantas vezes mais transistores poderão, em teoria, ser empacotados nessa nova pastilha de \textbf{2} nm.
\par\vspace{\espacoQuestao}

\end{multicols*}
\end{document}